\section{Experiments}
\label{sec:experiments}

\subsection{Experimental setup}

The controlled benchmark uses N-CMAPSS FD1 as the primary subset and FD2 as a secondary cross-subset check. Unless otherwise noted, the apples-to-apples rows use the same data pipeline described in Section~\ref{sec:methods}: last-window truncation, train-fit feature standardization, reader max RUL 65, Adam optimization~\cite{kingma2015adam}, and the shared callback stack. The final review-response suite uses three seeds (42--44), with metrics saved in the repository's review-response results directory. External comparison rows are cited only as contextual literature anchors, not as direct matched baselines.

\subsection{Architecture sweep}

Table~\ref{tab:families} summarizes the architecture families covered by the registry. The current codebase catalog contains 22 registered models, ranging from simple no-temporal baselines to multi-scale attention and transformer variants.

\begin{table}[t]
\centering
\caption{Architecture families in the benchmark harness.}
\label{tab:families}
\small
\resizebox{\linewidth}{!}{%
\begin{tabular}{lll}
\toprule
Family & Registered models & Role in this paper \\
\midrule
Baseline & MLP & No-temporal lower bound \\
RNN & LSTM~\cite{hochreiter1997lstm}, BiLSTM, GRU~\cite{cho2014gru}, BiGRU & Long-sequence recurrent baselines \\
Convolutional & TCN~\cite{bai2018tcn}, WaveNet~\cite{oord2016wavenet} & Dilated convolution baselines \\
Hybrid & CNN-LSTM, CNN-GRU, Inception-LSTM & Local feature extractor plus recurrence \\
Attention & Attention-LSTM, Transformer~\cite{vaswani2017attention}, ResNet-LSTM & Attention and transformer comparisons \\
Published fusion & MDFA, MDFA-paper, CNN-LSTM-Attention, CATA-TCN, TTSNet~\cite{ttsnet2025}, ATCN, MSTCN~\cite{xu2024msattn} & Multi-scale / attention architectures \\
Transformer variants & STAR~\cite{fan2024star}, sparse-transformer-BiGRCU & Recent hierarchical/sparse attention \\
\bottomrule
\end{tabular}
}
\end{table}

\begin{figure}[t]
\centering
\includegraphics[width=\linewidth]{architecture_leaderboard_fd1}
\caption{Historical FD1 architecture sweep from the repository comparison artifact. The figure is used as broad architecture-family context for recurrent~\cite{hochreiter1997lstm,cho2014gru}, convolutional~\cite{bai2018tcn,oord2016wavenet}, transformer~\cite{vaswani2017attention,fan2024star}, and multi-scale attention models~\cite{ttsnet2025,xu2024msattn}; the controlled head-to-head claims in the paper use the apples-to-apples rows in Tables~\ref{tab:fd1}--\ref{tab:fd2}.}
\label{fig:architecture-leaderboard}
\end{figure}

\subsection{FD1 headline results}

Table~\ref{tab:fd1} reports the three-seed FD1 head-to-head at $T{=}1000$. WaveNet is the strongest and most stable model in this comparison (RMSE $6.868\pm0.299$). MSTCN is stable and competitive (RMSE $7.890\pm0.526$), but it does not beat WaveNet on FD1. CNN-GRU has one strong seed and two collapsed seeds, so its mean is much worse than its best run.

\begin{table}[t]
\centering
\caption{Review-response FD1 results at $T{=}1000$ over seeds 42--44.
  All rows use the same preprocessing, reader max RUL 65, asymmetric MSE,
  and 30-epoch budget with early stopping. Std is over seeds.}
\label{tab:fd1}
\small
\begin{tabular}{llrrrr}
\toprule
Model & Seeds & RMSE mean & RMSE std & Mean $R^2$ & Acc@20 \\
\midrule
WaveNet & 3 & 6.868 & 0.299 & 0.8986 & 98.24 \\
MSTCN & 3 & 7.890 & 0.526 & 0.8659 & 97.75 \\
CNN-GRU & 3 & 17.168 & 8.801 & 0.2561 & 64.13 \\
\bottomrule
\end{tabular}
\end{table}

The FD1 result is therefore not a single-run leaderboard. WaveNet is robustly best among the three models tested here, MSTCN is a stable second cluster member, and CNN-GRU should not be described as a reliable top model despite its good seed-42 result.

\subsection{FD2 robustness}

On FD2, MSTCN is the best model by RMSE among the available apples-to-apples rows (T11: 6.495), with WaveNet close behind (T12: 6.739) and CNN-GRU degrading under the more variable flight envelope (T13: 19.812). WaveNet still has the highest Accuracy@20 on FD2 (T12: 99.50 versus T11: 99.01). Thus, the FD2 result supports the ``top cluster'' framing for MSTCN without establishing dominance across every metric.

\begin{table}[t]
\centering
\caption{Canonical FD2 cross-subset results.}
\label{tab:fd2}
\small
\begin{tabular}{llrrrr}
\toprule
Row & Model & RMSE & $R^2$ & PHM & Acc@20 \\
\midrule
T11 & MSTCN & 6.495 & 0.8897 & 0.5755 & 99.01 \\
T12 & WaveNet & 6.739 & 0.8813 & 0.6756 & 99.50 \\
T13 & CNN-GRU & 19.812 & -0.0265 & 5.1463 & 59.41 \\
\bottomrule
\end{tabular}
\end{table}

\begin{figure}[t]
\centering
\includegraphics[width=0.82\linewidth]{fd_rmse_comparison}
\caption{Controlled RMSE comparison on FD1 and FD2 for the three models available in both apples-to-apples runs. WaveNet is best on FD1, MSTCN is best by RMSE on FD2, and CNN-GRU degrades sharply on FD2.}
\label{fig:fd-rmse}
\end{figure}

\begin{figure}[t]
\centering
\includegraphics[width=\linewidth]{metric_tradeoff_acc_rmse}
\caption{Accuracy--RMSE tradeoff for the controlled FD1/FD2 rows. Marker size scales with normalized PHM score, so lower and smaller points are better on the asymmetric PHM metric.}
\label{fig:metric-tradeoff}
\end{figure}

\begin{figure}[t]
\centering
\includegraphics[width=0.88\linewidth]{phm_accuracy_profile}
\caption{Normalized PHM score and Accuracy@20 for the controlled rows. Most top-cluster rows are tightly grouped on Accuracy@20, while PHM exposes additional separation because it penalizes late predictions asymmetrically.}
\label{fig:phm-accuracy}
\end{figure}

\begin{figure}[t]
\centering
\includegraphics[width=0.88\linewidth]{runtime_rmse_tradeoff}
\caption{Runtime--RMSE tradeoff for the controlled FD1/FD2 rows. Runtime is measured by the benchmark script on the local machine and is reported as a practical engineering diagnostic, not as a hardware-normalized efficiency result.}
\label{fig:runtime-rmse}
\end{figure}

\subsection{Multiseed reproducibility}

The review-response suite was run specifically to bound run-to-run variance. Table~\ref{tab:fd1} shows that the ranking changes when seed variance is included: WaveNet remains strong, MSTCN remains stable, and CNN-GRU is not reproducible under the same settings. This directly addresses the concern that the original single-seed table overstated architecture rankings.

\begin{table}[t]
\centering
\caption{Legacy multiseed reproducibility figure source. The paper's final
  architecture ranking uses Table~\ref{tab:fd1}; this table is retained only
  to document the older seed-variance artifact used for Figure~\ref{fig:multiseed-rmse}.}
\label{tab:multiseed}
\small
\begin{tabular}{llrrrr}
\toprule
Row & Model & Seeds & Mean RMSE & Std RMSE & Mean $R^2$ \\
\midrule
T5 & MSTCN  & 3 (42--44) & 8.140 & 0.74 & 0.8565 \\
T6 & WaveNet & 2 (45--46) & 8.238 & 1.44 & 0.8498 \\
\bottomrule
\end{tabular}
\end{table}

\begin{figure}[t]
\centering
\includegraphics[width=0.72\linewidth]{multiseed_rmse_variability}
\caption{Per-seed RMSE for the FD1 multiseed runs. Dots show individual seeds and black squares show the mean with one standard deviation over the available seeds.}
\label{fig:multiseed-rmse}
\end{figure}

\subsection{Sequence-length evidence}

Table~\ref{tab:length} gives the clean sequence-length result requested by reviewers: one architecture (WaveNet), one training recipe, and three seeds. The sweep shows that short windows are effective across the tested range. The best mean is at $T{=}500$, with $T{=}250$ and $T{=}1000$ close enough that the result should be interpreted as a practical short-window regime rather than an exact optimum.

\begin{table}[t]
\centering
\caption{Matched WaveNet sequence-length sweep on FD1 over seeds 42--44.
  All rows use asymmetric MSE, 32 features, reader max RUL 65, and the same
  training budget. Std is over seeds.}
\label{tab:length}
\small
\begin{tabular}{lrrrr}
\toprule
Sequence length & RMSE mean & RMSE std & Mean $R^2$ & Acc@20 \\
\midrule
100 & 7.098 & 0.375 & 0.8916 & 97.95 \\
250 & 6.880 & 0.531 & 0.8980 & 98.44 \\
500 & 6.690 & 0.173 & 0.9038 & 98.24 \\
1000 & 6.868 & 0.299 & 0.8986 & 98.24 \\
\bottomrule
\end{tabular}
\end{table}

\begin{figure}[t]
\centering
\includegraphics[width=\linewidth]{sequence_length_evidence}
\caption{FD1 sequence-length evidence from the repository artifact. The final paper claim uses the matched WaveNet sweep in Table~\ref{tab:length}; older mixed-protocol rows are treated as historical context only.}
\label{fig:length-evidence}
\end{figure}

\subsection{External comparison}

The closest N-CMAPSS literature anchors are Solís-Martín et al.~\cite{solismartin2021stacked}, who report validation RMSE 6.24 using a 161-cycle window, and Cohen et al.~\cite{ncmapss_eventual2023}, who report RMSE 7.75 for ANN-Flux. Our three-seed FD1 WaveNet mean at $T{=}1000$ (RMSE $6.868\pm0.299$), the best matched window-sweep mean at $T{=}500$ (RMSE $6.690\pm0.173$), and the MSTCN FD1 mean (RMSE $7.890\pm0.526$) are in the same broad N-CMAPSS numerical band. These are not direct comparisons because external rows use their own splits, inputs, and validation/test conventions. C-MAPSS-only results from STAR~\cite{fan2024star}, TTSNet~\cite{ttsnet2025}, and Xu et al.~\cite{xu2024msattn} are treated as architectural precedent rather than direct numerical baselines because the datasets and sequence regimes differ substantially.

\begin{figure}[t]
\centering
\includegraphics[width=0.82\linewidth]{external_ncmapss_comparison}
\caption{Contextual comparison with the closest N-CMAPSS literature anchors~\cite{solismartin2021stacked,ncmapss_eventual2023}. External rows use their published protocols and are not direct head-to-head baselines; they are included only to show the numerical band of recent N-CMAPSS results.}
\label{fig:external-comparison}
\end{figure}

We dissect these results across sequence length, architecture family, and loss form in Section~\ref{sec:ablations}.
